\documentclass[12 pt, letterpaper]{hmcpset}
\usepackage{hw-preamble}
\usepackage{pdfpages}

\name{Jayden Thadani}
\class{MATH 189AA --- Commutative Algebra}
\assignment{Exam 2}
\duedate{25th November 2025}

\begin{document}

\begin{problem}[1.]
	Consider the ideal $I = (x, z)$ in the ring
	\[
		R = \mathbb{Q}[x, y, z] / (y^{2} - xz, x^{2} y - z^{2}, x^{3} - yz).
	\]
	\begin{enumerate}
		\item Use Macaulay2 to calculate a free resolution $P_{\bullet}$ of $I$ over $R$. Find the modules $P_{0}, P_{1}, P_{2}$ and the matrices representing the homomorphisms between them. Write your solution in the form
			\[
				\begin{tikzcd}
					\cdots \ar[r] & P_{2} \ar[r, "\text{matrix}"] & P_{1} \ar[r, "\text{matrix}"] & P_{0} \ar[r, "\text{matrix}"] & (x, z) \ar[r] & 0
				\end{tikzcd}
			\]

		\item Print your Macaulay2 session and hand it in with your exam.
	\end{enumerate}
\end{problem}

\begin{solution}
	\begin{enumerate}
		\item $P_{0} = R^{\oplus 2}$, $P_{1} = R^{\oplus 3}$, and $P_{2} = R^{\oplus 6}$. The morphism $P_{0} \to I$ is given by
			\[
				\begin{pmatrix}
					x & 0 \\
					0 & y
				\end{pmatrix}
			\]
			The morphism $P_{1} \to P_{0}$ is given by the matrix
			\[
				\begin{pmatrix}
					-z & - y & x^{2} \\
					y & x & -z
				\end{pmatrix}.
			\]
			The morphism $P_{2} \to P_{1}$ is given by
			\[
				\begin{pmatrix}
					y & -x & x^{2} & z & -z & 0 \\
					-z & y & 0 & 0 & x^{2} & z \\
					0 & 0 & z & y & 0 & x
				\end{pmatrix}
			\]

		\item My Macaulay2 session is included in the two pages below.
	\end{enumerate}
\end{solution}

\includepdf[pages = {1, 2}]{figures/P1 macaulay.pdf}

\pagebreak

\begin{problem}[2.]
	The diagram below is a commutative diagram of $R$-modules with exact rows:
	\[
		\begin{tikzcd}
			A \ar[r, "f_{1}"] \ar[d, "\alpha"] & B \ar[r, "f_{2}"] \ar[d, "\beta"] & C \ar[r, "f_{3}"] \ar[d, "\gamma"] & D \ar[r, "f_{4}"] \ar[d, "\delta"] & E \ar[d, "\varepsilon"] \\
			A' \ar[r, "f_{1}"] & B' \ar[r, "f_{2}"] & C' \ar[r, "f_{3}"] \ar[r] & D' \ar[r, "f_{4}"] & E'
		\end{tikzcd}
	\]
	Establish the following claims directly, by diagram chase:
	\begin{enumerate}
		\item If $\alpha$ is surjective and $\beta$ and $\delta$ are both injective, then $\gamma$ is injective.

		\item If $\varepsilon$ is injective and $\beta$ and $\delta$ are both surjective, then $\gamma$ is surjective.
	\end{enumerate}
\end{problem}

\pagebreak

\begin{problem}[3.]
	If $P$ is a finitely generated projective $R$-module, show that $\operatorname{Hom}_{R}(P, R)$ is also a projective $R$-module.
\end{problem}

\begin{solution}
	Since $P$ is projective, we can find some $R$-module $N$ such that $P \oplus N \cong F$ is free. Thus,
	\[
		\operatorname{Hom}_{R}(F, R) \cong \operatorname{Hom}_{R}(P, R) \times \operatorname{Hom}_{R}(N, R) \cong \operatorname{Hom}_{R}(P, R) \oplus \operatorname{Hom}_{R}(N, R).
	\]
	Further, since $P$ is finitely generated, we can choose $F$ to be free. We claim that this makes $\operatorname{Hom}_{R}(F, R)$ free. Then, since we have written $\operatorname{Hom}_{R}(P, R)$ as a direct summand of a free module, we have shown that it is projective.

	The freeness of $\operatorname{Hom}_{R}(F, R)$ follows by the chain of isomorphisms below.
	\begin{align*}
		\operatorname{Hom}_{R}(F, R) & \cong \operatorname{Hom}_{R}(R^{\oplus n}, R) \\
					     & = \bigoplus_{i = 1}^{n} \operatorname{Hom}_{R}(R, R) \\
					     & = R^{\oplus n}.
	\end{align*}
	The last equality follows because $\operatorname{Hom}_{R}(R, M) \cong M$.
\end{solution}

\pagebreak

\begin{problem}[4.]
	Let
	\[
		\begin{tikzcd}
			0 \ar[r] & L \ar[r] & M \ar[r] & N \ar[r] & 0
		\end{tikzcd}
	\]
	be a short exact sequence of $R$-modules. If $M$ and $N$ are projective, prove that $L$ is also projective.
\end{problem}

\begin{solution}
	Since $N$ is projective, the exact sequence splits. Thus, we have $M \cong L \oplus N$. Since $M$ is projective, it is a direct summand of some free module $F$. But then, as a direct summand of $M$, $L$ also becomes a direct summand of $F$ (if $M \oplus Q \cong F$, then $L \oplus (N \oplus Q) \cong F$). That is, $L$ is projective.
\end{solution}

\pagebreak

\begin{problem}[5.]
	Suppose $M$ is a finitely presented $R$-module and let
	\[
		\begin{tikzcd}
			0 \ar[r] & K \ar[r] & N \ar[r] & M \ar[r] & 0
		\end{tikzcd}
	\]
	be a short exact sequence where $N$ is finitely generated. Prove that $K$ is also finitely generated.
\end{problem}

\begin{solution}
	Denote the map $N \to M$ by $\varphi$. Notice that then $K = \ker \varphi$.

	Since $M$ is finitely presented, there is an exact sequence
	\[
		\begin{tikzcd}
			R^{\oplus n} \ar[r, "\theta"] & R^{\oplus m} \ar[r, "\psi"] & M \ar[r] & 0.
		\end{tikzcd}
	\]
	We claim that we can construct maps $\alpha$ and $\beta$ so that the following diagram commutes.
	\[
		\begin{tikzcd}
			& R^{\oplus n} \ar[r, "\theta"] \ar[d, dashed, "\alpha"] & R^{\oplus m} \ar[r, "\psi"] \ar[d, dashed, "\beta"] & M \ar[r] \ar[d, "\operatorname{id}_{M}"] & 0 \\
			0 \ar[r] & K \ar[r, "i"] & N \ar[r, "\varphi"] & M \ar[r] & 0
		\end{tikzcd}
	\]
	$\beta$ exists by the projectivity of $R^{\oplus m}$ --- it is the lift of $\psi : R^{\oplus m} \to M$ to a map $R^{\oplus m} \to N$ (using the surjectivity of $\varphi : N \to M$). Note that this property makes the right square commute.

	To show the existence of $\alpha$ we use the universal property of the kernel of $\varphi : M \to N$. That is, any map $\pi : L \to N$ such that $\varphi \circ \pi : L \to N \to M$ is the zero map must factor through the kernel map $i : K \to N$.

	We claim that $\beta \circ \theta$ is such a map. Specifically,
	\begin{align*}
		\varphi \circ (\beta \circ \theta) & = (\varphi \circ \beta) \circ \theta \\
						   & = \operatorname{id}_{M} \circ \psi \circ \theta \\
						   & = \operatorname{id}_{M} \circ 0 \\
						   & = 0.
	\end{align*}
	Here, the second equality follows by the commutativity of the right square. Thus, by the universal property of the kernel, there is a unique $\alpha : R^{\oplus m} \to K$ such that $i \circ \alpha = (\beta \circ \theta)$. This is exactly a map $\alpha$ so that the left square commutes.

	Then applying the snake lemma, we get an exact sequence
	\[
		\begin{tikzcd}
			\ker \operatorname{id}_{M} = 0 \ar[r] & \operatorname{coker} \alpha \ar[r] & \operatorname{coker} \beta \ar[r] & 0 = \operatorname{coker} \operatorname{id}_{M}
		\end{tikzcd}
	\]
	That is, $\operatorname{coker} \alpha \cong \operatorname{coker} \beta$. Since the finitely generated module $N$ surjects onto $\operatorname{coker} \beta$, it and $\operatorname{coker} \alpha$ are finitely generated.

	Finally, we have $K / \operatorname{img} \alpha \cong \operatorname{coker} \alpha$. Note that $\operatorname{img} \alpha$ is finitely generated since $R^{\oplus m}$ surjects onto it. Thus, the quotient of $K$ by a finitely generated submodule is finitely generated. This implies $K$ is finitely generated.
\end{solution}

\end{document}
